\section{Clase 9 - Pr\'actica 4 (N\'umero de condici\'on)}

\subsection{M\'etodos directos para la resoluci\'on de sistemas lineales}

\underline{En general}:

$$ L_{1} =
	\begin{pmatrix}
        1 & \cdots & \cdots & \cdots & \cdots\\
        -l_{21} & \cdots & \cdots & \cdots & \cdots\\
        \vdots & \vdots & \vdots & \vdots & \vdots\\
        -l_{n1} & \cdots & \cdots &\cdots & 1\\
	\end{pmatrix}
	,\ \text{con}\ l_{i1} = \frac{a_{i1}}{a_{11}}\ \text{si}\ a_{11} \neq 0
$$

$$ L_{2} =
	\begin{pmatrix}
        1 & \cdots & \cdots & \cdots & \cdots\\
        \vdots & \vdots & \vdots & \vdots & \vdots\\
        \cdots & -l_{32} & \cdots & \cdots & \cdots\\
        \vdots & \vdots & \vdots & \vdots & \vdots\\
        \cdots & -l_{n2} & \cdots & \cdots & 1\\
	\end{pmatrix}
	,\ \text{con}\ l_{i2} = \frac{a_{i2}^{(2)}}{a_{22}^{(2)}}\ \text{si}\ a_{22}^{(2)} \neq 0
$$

Notemos que $a_{22}^{(2)} = -l_{21}a_{12} + a_{22} \neq 0 \iff a_{11}a_{22} - a_{21}a_{12} \neq 0$.\\

Si $ A_{k}^{(k)} =
	\begin{pmatrix}
        a_{1k}^{(k)}\\
        \vdots\\
        a_{kk}^{(k)}\\
        a_{k + 1, k}^{(k)}\\
        a_{nk}^{(k)}\\
	\end{pmatrix}
$, buscamos $L_{k}$ tal que $ L_{k}A_{k}^{(k)} =
	\begin{pmatrix}
        a_{1k}^{(k + 1)}\\
        \vdots\\
        a_{kk}^{(k + 1)}\\
        0\\
        0\\
	\end{pmatrix}
$. Luego $l_{ik} = \frac{a_{ik}^{(k)}}{a_{kk}^{(k)}}$ para $k \leq i \leq n$.

$$ L_{k} =
	\begin{pmatrix}
        1 & \cdots & \cdots & \cdots & \cdots\\
        \vdots & 1 & \vdots & \vdots & \vdots\\
        \cdots & -l_{k + 1, k} & \cdots & \cdots & \cdots\\
        \vdots & \vdots & \vdots & \vdots & \vdots\\
        \cdots & -l_{nk} & \cdots & \cdots & 1\\
	\end{pmatrix}
	=
	I -
	\begin{pmatrix}
        0\\
        \vdots\\
        0\\
        l_{k + 1, k}\\
        l_{nk}\\
	\end{pmatrix}
	\begin{pmatrix}
        0, & \cdots, & 1, & \cdots, & 0\\
	\end{pmatrix}
	= I - l_{k}e_{k}^{t}
$$

$$(I - l_{k}e_{k}^{t})(I + l_{k}e_{k}^{t}) = I - l_{k}e_{k}^{t}l_{k}e_{k}^{t} = I \Rightarrow L_{k}^{t} = I + l_{k}e_{k}^{t}$$

Adem\'as, $L_{k}^{-1}L_{k + 1}^{-1} = (I + l_{k}e_{k}^{t})(I + l_{k + 1}e_{k + 1}^{t}) = I + l_{k + 1}e_{k + 1}^{t} + l_{k}e_{k}^{t} + l_{k}e_{k}^{t}l_{k + 1}e_{k + 1}^{t}$, entonces:

$$ L_{k}^{-1}L_{k + 1}^{-1} =
	\begin{pmatrix}
        1 & \cdots & \cdots & \cdots & \cdots\\
        \vdots & 1 & \vdots & \vdots & \vdots\\
        \cdots & l_{k + 1, k} & 1 & \cdots & \cdots\\
        \vdots & \vdots & l_{k + 1, k + 1} & \vdots & \vdots\\
        \cdots & l_{nk} & l_{n,k + 1} & \cdots & 1\\
	\end{pmatrix}
$$

Luego,

$$ L =
	\begin{pmatrix}
        1 & \cdots & \cdots & \cdots & \cdots\\
        l_{21} & \cdots & \cdots & \cdots & \cdots\\
        \vdots & l_{32} & \vdots & \vdots & \vdots\\
        l_{n1} & l_{n2} & \cdots & l_{n, n - 1} & 1\\
	\end{pmatrix}
$$\\

\underline{Algoritmo de eliminaci\'on gaussiana sin pivoteo}\\

\begin{verbatim}
    U = A, L = I

    para k = 1 hasta n - 1:
       para i = k + 1 hasta n:
           L[i, k] = U[i, k] / U[k, k]
           U[i, k:n] = U[i, k:n] - L[i, k] * U[k, k:n]
\end{verbatim}

Son aproximadamente $\frac{2}{3}n^{3}$ operaciones (Trefethen \& Bau, p. 152).\\

\underline{Teorema}: Sea $A \in \mathbb{R}^{n \times n}$ tal que en el procedimiento de eliminaci\'on gaussiana no aparecen pivotes nulos ($a_{kk}^{(k)} \neq 0,\ k = 1, \cdots, n$). Entonces existen matrices \'unicas $L$ y $U$, $L$ triangular inferior con unos en la diagonal y $U$ triangular superior tal que $A = LU$.\\

\underline{Demostraci\'on}: Veamos la unicidad: $\sup A = LU = \tilde{L}\tilde{U}$:

$$\text{det}(A) = \text{det}(L) = \text{det}(U) = \text{det}(\tilde{L})\text{det}(\tilde{U}) = \prod \limits_{k = 1}^{n} a_{kk}^{(k)} \neq 0$$

Luego, $U$ y $\tilde{U}$ (y $L$ y $\tilde{L}$) son inversibles. Tenemos:

$$\tilde{L}^{-1}LUU^{-1} = \tilde{L}^{-1}\tilde{L}\tilde{U}U^{-1} \Rightarrow \tilde{L}^{-1}L = D = \tilde{U}U^{-1} \Rightarrow D = I \Rightarrow L = \tilde{L}\ \text{y}\ U = \tilde{U}$$

con $\tilde{L}^{-1}L$ matriz triangular inferior con unos en la diagonal, $D$ diagonal y $\tilde{U}U^{-1}$ triangular superior.\\

Si $ A =
	\begin{pmatrix}
        0 & 2 & 1\\
        3 & -4 & 1\\
        -1 & 1 & 1\\
	\end{pmatrix}
$, como $a_{11} = 0$, no puedo hacer eliminaci\'on gaussiana. Permuto filas $P = P_{12}$ y obtengo $ PA =
	\begin{pmatrix}
        3 & -4 & 1\\
        0 & 2 & 1\\
        -1 & 1 & 1\\
	\end{pmatrix}
$, y ahora s\'i.\\

\underline{Teorema}: Sea $A \in \mathbb{R}^{n \times n}$ inversible, existen una matriz de permutaci\'on $P$, una matriz $L$ triangular inferior con unos en la diagonal y una matriz $U$ triangular superior tales que $PA = LU$.\\

\textit{Ejemplo con permutaci\'on}:

$$
	\begin{pmatrix}
        1 & 2 & 0\\
        2 & 4 & 3\\
        0 & 5 & 6\\
	\end{pmatrix}
	\rightarrow_{
	\tilde{F}_{2} = F_{2} - 2F_{1}}
    \begin{pmatrix}
        1 & 2 & 0\\
        0 & 0 & 3\\
        0 & 5 & 6\\
	\end{pmatrix}
$$

No podemos seguir con la permutaci\'on. Permutamos fila 2 y fila 3:

$$
	\begin{pmatrix}
        1 & 2 & 0\\
        2 & 4 & 3\\
        0 & 5 & 6\\
	\end{pmatrix}
	\rightarrow_{
	F_{2} \leftrightarrow F_{3}}
    \begin{pmatrix}
        1 & 2 & 0\\
        0 & 5 & 6\\
        2 & 4 & 3\\
	\end{pmatrix}
	\rightarrow_{
	\tilde{\tilde{F}}_{2} = \tilde{F}_{3} - 2\tilde{F}_{1}}
	\begin{pmatrix}
        1 & 2 & 0\\
        0 & 5 & 6\\
        0 & 0 & 3\\
	\end{pmatrix} = U
$$

$$ P_{1} = 
	\begin{pmatrix}
        1 & 2 & 0\\
        0 & 0 & 1\\
        0 & 1 & 0\\
	\end{pmatrix}
	,\ L_{1} = 
    \begin{pmatrix}
        1 & 0 & 0\\
        0 & 1 & 0\\
        -2 & 0 & 1\\
	\end{pmatrix}
	\Rightarrow L_{1}P_{1}A = U \Rightarrow P_{1}A = L_{1}^{-1}U \Rightarrow PA = LU,\ \text{con}\ P = P_{1}\ \text{y}\ L = L_{1}^{-1}
$$

Si tenemos por ejemplo:

$$L_{3}P_{3}L_{2}P_{2}L_{1}P_{1}A = U$$

podemos reescribir::

$$L_{3}P_{3}L_{2}P_{2}L_{1}P_{1} = L_{3}'L_{2}'L_{1}P_{3}'P_{2}'P_{1}'$$

con $L'_{3} = L_{3},\ L'_{2} = P_{3}L_{2}P_{3}^{-1},\ L_{1}' = P_{3}P_{2}L_{1}P_{2}^{-1}P_{3}^{-1}$. Y as\'i:

$$P_{3}P_{2}P_{1}A = (L'_{3}L'_{2}L'_{1})^{-1}U \Rightarrow PA = LU$$

con $P = P_{3}P_{2}P_{1}$ y $ L = (L'_{3}L'_{2}L'_{1})^{-1}$.

\subsubsection{Cholesky}

\underline{Definici\'on}: $A$ se dice definida positiva si $x^{t}Ax = \langle x, Ax \rangle > 0\ \forall \ x \neq 0$.\\

\underline{Observaci\'on}: Si $A = LL^{t}$ es inversible, entonces $A$ es sim\'etrica y definida positiva:

$$(LL^{t})^{t} = (L^{t})^{t}L^{t} = LL^{t}$$

$$\langle x, LL^{t}x \rangle = \langle L^{t}x, L^{t}x \rangle = \|L^{t}x\|^{2} > 0$$

\underline{Teorema (Cholesky)}: Sea $A \in \mathbb{R}^{n \times n}$ sim\'etrica y definida positiva, entonces existe una \'unica matriz $L$ triangular inferior con $l_{ii} > 0$ para $i = 1, \cdots, n\ / \ A = LL^{t}$.\\

\underline{Teorema 1}: Si $A \in \mathbb{R}^{n \times n}$. Sea $ M_{k} =
	\begin{pmatrix}
        a_{11} & \cdots & a_{1k}\\
        \vdots & \vdots & \vdots\\
        a_{k1} & \cdots & a_{kk}\\
	\end{pmatrix}
$. Si $\text{det}(M_{k}) \neq 0$ para $k = 1, \cdots, n$ entonces todos los pivotes de $A$ son no nulos ($a_{kk}^{(k)} \neq 0$ para $k = 1, \cdots, n$).\\

\underline{Demostraci\'on}: $\text{det}(M_{1}) = a_{11} \neq 0$. Supongamos que se puede construir $L_{1}, \cdots, L_{k - 1}$, queremos ver que puedo construir $L_{k}$ (con $a_{kk}^{(k)} \neq 0$):

$$M_{k} =
	\begin{pmatrix}
        a_{11} & \cdots & a_{1k}\\
        \vdots & \vdots & \vdots\\
        a_{k1} & \cdots & a_{kk}\\
	\end{pmatrix}
	=
	\begin{pmatrix}
        1 & \cdots & \cdots & \cdots\\
        l_{21} & \cdots & \cdots & \cdots\\
        \vdots & \vdots & \vdots & \vdots\\
        l_{k1} & \cdots & l_{k, k - 1} & 1\\
	\end{pmatrix}
	\begin{pmatrix}
        a_{11} & \cdots & \cdots & a_{1k}\\
        0 & a_{22}^{2} & \cdots & a_{2k}^{(2)}\\
        \vdots & \vdots & \vdots & \vdots\\
        0 & \cdots & 0 & a_{kk}^{(k)}\\
	\end{pmatrix}
$$

Luego, $\text{det}(M_{k}) = \prod \limits_{j = 1}^{k} a_{jj}^{(j)}$ y deducimos que $a_{kk}^{(k)} \neq 0$.\\

\underline{Teorema 2}: Sea $A$ definida positiva, entonces $M_{k}$ es definida positiva ($1 \leq k \leq n$).\\

\underline{Demostraci\'on}: $x \in \mathbb{R}^{k} - \{\Vec{0}\},\ x^{t}M_{k}x = \begin{pmatrix}
        x^{t}, & \Vec{0}\\
       	\end{pmatrix}A\begin{pmatrix}
        x\\
        \Vec{0}\\
       	\end{pmatrix} > 0$.\\

\underline{Lema 3}: Si $A$ es definida positiva, entonces $a_{jj} > 0 (1 \leq j \leq n)$.\\

\underline{Demostraci\'on}: $0 < e_{j}^{t}Ae_{j} = a_{jj}$.\\

\underline{Demostraci\'on (Cholesky)}: Por los teoremas 1 y 2, como los pivotes son no nulos, existen matrices \'unicas $L$ y $U$, $L$ triangular inferior con unos en la diagonal y $U$ triangular superior tales que $A = LU$.\\

Como $A$ es sim\'etrica, $A = LU = A^{t} = U^{t}L^{t}$ y por lo tanto:

$$LU = U^{t}L^{t} \Rightarrow L^{-1}LU(L^{t})^{-1} = L^{-1}U^{t}L^{t}(L^{t})^{-1} \Rightarrow L^{-1}U^{t} = U(L^{t})^{-1} = D$$

$$\Rightarrow U = DL^{t} \Rightarrow A = LDL^{t}$$

con $L^{-1}U^{t}$ triangular inferior, $U(L^{t})^{-1}$ triangular superior y $D$ diagonal.\\

Veamos que $D$ es definida positiva: sea $x \in \mathbb{R}^{n} - \{\Vec{0}\}$, como $L^{t}$ es inversible, $\exists\ \tilde{x} \in \mathbb{R}^{n} - \{\Vec{0}\}\ / \ L^{t}\tilde{x} = x \Rightarrow x^{t}Dx = \tilde{x}^{t}LDL^{t}\tilde{x} = \tilde{x}^{t}A\tilde{x} > 0$.\\

Luego por Lema 3, $d_{jj} > 0 (1 \leq j \leq n)$. Sea $ D^{\frac{1}{2}} =
	\begin{pmatrix}
        \sqrt{a_{11}} & \cdots & \cdots\\
        \vdots & \vdots & \vdots\\
        \cdots & \cdots & \sqrt{a_{nn}}\\
	\end{pmatrix}
\Rightarrow A = LDL^{t} = LD^{\frac{1}{2}}D^{\frac{1}{2}}L^{t}$.\\

Si $\tilde{L} = LD^{\frac{1}{2}} \Rightarrow A = \tilde{L}\tilde{L}^{t}$ y vale que $\tilde{L}_{ii} = \sqrt{d_{ii}} > 0$ para $1 \leq i \leq n$.\\

Veamos la unicidad: supongamos que existen $\tilde{L}_{1}$ y $\tilde{L}_{2}$ con $(\tilde{L}_{1})_{ii} > 0$, $(\tilde{L}_{2})_{ii} > 0$ para $1 \leq i \leq n$, tales que $A = \tilde{L}_{1}\tilde{L}_{1}^{t} = \tilde{L}_{2}\tilde{L}_{2}^{t}$:

$$\tilde{L}_{2}^{-1}\tilde{L}_{1} = \tilde{L}_{2}^{t}(\tilde{L}_{1}^{t})^{-1} = \tilde{D}$$

con $\tilde{L}_{2}^{-1}\tilde{L}_{1}$ triangular inferior, $\tilde{L}_{2}^{t}(\tilde{L}_{1}^{t})^{-1}$ triangular superior y $\tilde{D}$ diagonal con $d_{ii} > 0$ para $1 \leq i \leq n$.

$$\tilde{D} = \tilde{L}_{2}^{-1}\tilde{L}_{1} = (\tilde{L}_{1}^{-1}\tilde{L}_{2})^{-1} = ((\tilde{L}_{2}^{t}(\tilde{L}_{1}^{t})^{-1})^{t})^{-1} = (\tilde{D}^{t})^{-1} = \tilde{D}^{-1}$$

y por lo tanto, $\tilde{D} = I$ y $\tilde{L}_{1} = \tilde{L}_{2}$.\\

\underline{Algoritmo de Cholesky}

\begin{verbatim}
    L = I
    L[1, 1] = sqrt(A[1, 1])

    para t = 2 hasta n:
        # llenamos la primer columna
        L[t, 1] = A[t, 1] / L[1, 1]
        
    para j = 2 hasta n:
        # llenamos las otras columnas de arriba hacia abajo
        L[j, j] = sqrt(A[j, j] - sum((L[j, 1:j - 1])**2))
        
        para i = j + 1 hasta n:
            L[i, j] = (A[i, j] - sum(L[i, 1: j - 1].L[j, 1: j - 1])) / L[j, j]
\end{verbatim}

\underline{Observaci\'on}: El n\'umero de operaciones es $\frac{n^{3}}{6}$ (la mitad de Gauss).\\

\textit{Ejemplo}:

$$A =
	\begin{pmatrix}
        4 & 12 & -16\\
        12 & 37 & -43\\
        -16 & -43 & 98\\
	\end{pmatrix}
$$

\begin{itemize}
	\item Sim\'etrica \checkmark
	
	\item Definida positiva: $A$ es definida positiva $\iff$ todos sus menores principales son positivos
	
	$$4\ \checkmark$$
	$$4\cdot37 - 12\cdot12 = 4 > 0\ \checkmark$$
	$$\text{det}(A) = 36 > 0\ \checkmark$$
\end{itemize}

Busquemos la descomposici\'on:

$$A =
	\begin{pmatrix}
        4 & 12 & -16\\
        12 & 37 & -43\\
        -16 & -43 & 98\\
	\end{pmatrix}
	=
	\begin{pmatrix}
        l_{11} & 0 & 0\\
        l_{21} & l_{22} & 0\\
        l_{31} & l_{32} & l_{33}\\
	\end{pmatrix}
	\begin{pmatrix}
        l_{11} & l_{21} & l_{31}\\
        0 & l_{22} & l_{32}\\
        0 & 0 & l_{33}\\
	\end{pmatrix}
	= LL^{t}
$$

$$ 
	\begin{cases}
		l_{11}^{2} = 4 \Rightarrow l_{11} = 2\\
		l_{11}l_{21} = 12 \Rightarrow l_{21} = 6\\
		l_{11}l_{31} = -16 \Rightarrow l_{31} = -8\\
		36 + l_{22} = 37 \Rightarrow l_{22} = 1\\
		-48 + l_{22}l_{32} = -43 \Rightarrow l_{22}l_{32} = 5 \Rightarrow l_{32} = 5\\
		64 + l_{32}^{2} + l_{33}^{2} = 98 \Rightarrow 5^{2} + l_{33}^{2} = 34 \Rightarrow l_{33}^{2} = 9 \Rightarrow l_{33} = 3\\
    \end{cases}
$$

Luego,

$$L =
	\begin{pmatrix}
        2 & 0 & 0\\
        6 & 1 & 0\\
        -8 & 5 & 3\\
	\end{pmatrix}
$$